\section{Συνδυασμός \en{LLMs}, \en{LoRA} και \en{Federated Learning}}

\subsection{Εισαγωγή}

Η ενοποίηση των Μεγάλων Γλωσσικών Μοντέλων (\en{LLMs}), τεχνικών αποδοτικής προσαρμογής παραμέτρων όπως το \en{Low-Rank Adaptation (LoRA)} και του \en{Federated Learning (FL)} αποτελεί μια πολλά υποσχόμενη κατεύθυνση για την ανάπτυξη επεκτάσιμων και φιλικών προς την ιδιωτικότητα ευφυών συστημάτων. Τα \en{LLMs} προσφέρουν ισχυρές δυνατότητες επεξεργασίας μη δομημένων δεδομένων, όπως είναι τα αρχεία καταγραφής δικτύου, αλλά συνοδεύονται από υψηλές απαιτήσεις σε υπολογιστικούς πόρους και επικοινωνία. Το \en{Federated Learning}, από την άλλη πλευρά, επιτρέπει συνεργατική εκπαίδευση χωρίς ανταλλαγή ακατέργαστων δεδομένων, καθιστώντας το κατάλληλο για τομείς όπως η κυβερνοασφάλεια. Ωστόσο, απαιτεί αποδοτικές ενημερώσεις μοντέλων ώστε να μειώνεται το κόστος επικοινωνίας \cite{mcmahan2017communication}, γεγονός που δημιουργεί μια φυσική συνέργεια με τεχνικές \en{PEFT}.

\subsection{Κίνητρο}

Τα \en{LLMs} είναι ιδιαίτερα απαιτητικά υπολογιστικά λόγω του μεγάλου αριθμού παραμέτρων που διαθέτουν, με αποτέλεσμα η πλήρης προσαρμογή τους σε περιβάλλον \en{Federated Learning} να οδηγεί σε υπερβολικό κόστος επικοινωνίας και σημαντικές απαιτήσεις σε μνήμη και επεξεργαστική ισχύ. Σε περιβάλλοντα \en{5G} και \en{6G}, όπου πολλές συσκευές διαθέτουν περιορισμένους πόρους και εύρος ζώνης, απαιτούνται αποδοτικοί μηχανισμοί εκπαίδευσης ώστε η κατανεμημένη χρήση τέτοιων μοντέλων να είναι πρακτικά εφικτή.

\subsection{Ρόλος του \en{LoRA} στο \en{Federated Learning}}

Το \en{LoRA} μειώνει δραστικά τον αριθμό των εκπαιδεύσιμων παραμέτρων, καθώς αντί για πλήρεις ενημερώσεις του μοντέλου, οι πελάτες μεταδίδουν μόνο τους πίνακες χαμηλού βαθμού που εισάγονται για την προσαρμογή \cite{hu2021lora}. Η προσέγγιση αυτή οδηγεί σε σημαντικά μειωμένο κόστος επικοινωνίας, χαμηλότερες απαιτήσεις μνήμης και ταχύτερη εκπαίδευση και συνάθροιση ενημερώσεων. Με τον τρόπο αυτό καθίσταται δυνατή η αξιοποίηση \en{LLMs} σε κατανεμημένα περιβάλλοντα με περιορισμένους πόρους.

\subsection{Ενσωμάτωση σε Επίπεδο Συστήματος}

Σε ένα συνδυαστικό σύστημα, κάθε πελάτης διατηρεί ένα τοπικό αντίγραφο του προεκπαιδευμένου \en{LLM} και εφαρμόζει προσαρμογή μέσω \en{LoRA} χρησιμοποιώντας ιδιωτικά δεδομένα. Μόνο οι παράμετροι του \en{LoRA} κοινοποιούνται στον κεντρικό εξυπηρετητή, ενώ το βασικό μοντέλο και τα δεδομένα παραμένουν τοπικά. Η διαδικασία περιλαμβάνει τη διανομή του καθολικού μοντέλου, την τοπική εκπαίδευση με \en{LoRA} και τη συνάθροιση των ενημερώσεων μέσω \en{FL}. Η αρχιτεκτονική αυτή επιτρέπει κατανεμημένη προσαρμογή μοντέλων με διατήρηση της ιδιωτικότητας και αποδοτική χρήση πόρων, καθιστώντας την ιδιαίτερα κατάλληλη για ανάλυση αρχείων καταγραφής σε μεγάλης κλίμακας τηλεπικοινωνιακά δίκτυα.
